\documentclass[10pt,a4paper]{amsart}
\usepackage[margin=19mm]{geometry}
\usepackage{amsmath,amssymb,mathtools}
\usepackage[scr=rsfs,cal=euler]{mathalfa}
\usepackage{xfrac}
\usepackage[unicode,hidelinks,bookmarks=false]{hyperref}
\newcommand{\ie}{i.e.\ }
\newcommand{\cf}{cf.\ }
\newcommand{\resp}{respectively\ }
\newcommand{\Subsection}{\subsection}
\providecommand{\nobreakdash}{\nobreak\hbox{-}}
\setlength{\emergencystretch}{2em}
\multlinegap0pt
\usepackage{bm}
\usepackage{bbm}
\usepackage{dsfont}
\usepackage{xspace}
\usepackage{cancel}
\usepackage{mathtools}
\usepackage{amsthm}
\makeatletter
\@ifundefined{Theorem}{\newtheorem{Theorem}{Theorem}}{}
\@ifundefined{Lemma}{\newtheorem{Lemma}[Theorem]{Lemma}}{}
\makeatother
\title[Finiteness of Totally Magnetic Hypersurfaces]{Finiteness of Totally Magnetic Hypersurfaces}
\author{James Marshall Reber, Ivo Terek}
\date{Source: arXiv:2602.00439v2 (CC BY 4.0)}
\begin{document}
\maketitle

\noindent\emph{Source and licence.} Finiteness of Totally Magnetic Hypersurfaces. Version arXiv:2602.00439v2 (CC BY 4.0), distributed under the Creative Commons Attribution 4.0 International licence, as recorded in the arXiv licence metadata for this exact version and shown on the abstract page https://arxiv.org/abs/2602.00439v2. Full rights notice: licenses/arxiv-2602.00439-CC-BY-4.0.txt.\par\smallskip

\noindent\textit{Editorial scope.} Counted unit: the Lemma whose source label is \texttt{lem:dynamical-exp} in the article (source0.tex lines 328--330, source label \texttt{lem:dynamical-exp}), delivered together with the article's own complete original proof of it (source0.tex lines 332--357, one \texttt{proof} environment of 26 lines). No mathematics is written, repaired or re-derived: statement and proof are the article's own text line for line. Required context retained uncounted: Dc'/dt=XV (lines 261--263); def-psi-exp (lines 324--324). Every definition, display and label of those lines is retained unchanged. Omitted from this package and not redistributed: 1--260, 264--323, 325--327, 331--331, 358--573. Nothing inside a retained excerpt is omitted or altered: every retained body is delivered exactly as the article's lines are written, so no omission receipt of any kind accompanies this package. Citations: the retained excerpts carry no citation command at all, so no bibliography entry of the article is transcribed and no reference is redistributed. Reference adaptation: none was needed and none was made; every reference inside the delivered unit resolves inside it, so no unresolved reference or citation is produced. Printed slips kept literal: no printed word, symbol, hypothesis or number of the counted statement or of its proof was altered. Layout and class: the article's own class is \texttt{sigma} with packages including ifpdf, enumerate; the wrapper is the lane's pinned \texttt{amsart} editorial wrapper at 10pt on A4, which restates the theorem-like environments the retained text needs and adds the article's own macro definitions that the retained text needs (none), with extra wrapper packages (bm, bbm, dsfont, xspace, cancel, mathtools). The delivered numbering is the unit's own. Assets: no retained span contains a figure environment, a graphics-inclusion command, a PDF-page inclusion, a table or a TikZ picture; no image file is copied into this package, the package declares no asset dependency and no image file is redistributed with it. Licence: version arXiv:2602.00439v2 of this article is distributed under the Creative Commons Attribution 4.0 International licence, as recorded in the arXiv licence metadata for this exact version and shown on the abstract page https://arxiv.org/abs/2602.00439v2. A full rights notice is delivered at licenses/arxiv-2602.00439-CC-BY-4.0.txt. Characters: every retained excerpt is pure ASCII, so the delivered engine, XeLaTeX with its default Latin Modern text font, reports no missing character.
 \begin{equation}\label{Dc'/dt=XV}
 \frac{{\rm D}\dot{\gamma}}{{\rm d}t}(t)= X_V(\gamma(t),\dot{\gamma}(t)).
 \end{equation}

\begin{equation}\label{def-psi-exp} \exp_x^{\varphi}(v) = \begin{cases} \pi \circ \varphi^{\|v\|}(x,v/\|v\|) &\text{if } v \neq 0, \\ x & \text{if } v = 0. \end{cases}\end{equation}

\begin{Lemma} \label{lem:dynamical-exp}
 Assume that $\varphi^t$ is of class $C^k$, with $k \geq 2$ $($or $C^\infty$ or $C^\omega)$. Then for any point $x\in M$, $\exp^\varphi_x$ has the same regularity as $\varphi^t$ on $T_xM\smallsetminus \{0\}$, and it is of class $C^1$ on $T_xM$. Moreover, if it is of class $C^2$ at $0$ and $\varphi^t$ is an odd semi-spray flow, then $X_V(x,v) = 0$ for all $v\in U_xM$.
\end{Lemma}

\begin{proof}
 The $C^k$-regularity of $\exp^{\varphi}_x$ on $T_xM \smallsetminus \{0\}$ follows from \eqref{def-psi-exp}. In addition, all of its directional derivatives at $0$ exist as
 \begin{equation}\label{d0-psi-exp} {\rm d}_0\exp_x^{\varphi}(v) = \frac{\rm d}{{\rm d}t} \Big|_{t=0} \pi \circ \varphi^t(x,v) = X_H(x,v)\end{equation}for all $(x,v)\in UM$. Fixing coordinates $x^1, \ldots, x^n$ for $M$ centered at $x$, set
 \begin{equation}\label{gamma-i}
 \gamma^i(t,x,v) = \exp_x^{\varphi, i}(tv),
 \end{equation}where $\exp_x^{\varphi, i}(tv) = x^i(\exp_x^{\varphi}(tv))$; note that each $\gamma^i$ is of class $C^k$. Relative to linear coordinates $v^1,\ldots, v^n$ on $T_xM$, we may differentiate \eqref{gamma-i} with respect to $v^j$ to obtain
 \begin{equation}\label{d-gamma-i-d-vj}
 \frac{\partial \gamma^i}{\partial v^j}(t,x,v) =t \frac{\partial \exp_x^{\varphi,i}}{\partial v^j}(tv) \qquad\mbox{for all }t>0.
 \end{equation}Observe, however, that \eqref{d-gamma-i-d-vj} also holds at $t=0$, as $\gamma^i(0,x,v) = 0$, which is another consequence of \eqref{def-psi-exp} together with our choice of coordinates. With $(\partial \gamma^i/\partial t)(0,x,v) = X_H^i(x,v)$ for all $(x,v)\in UM$, where $X_H^i(x,v)$ denotes the $i$-th component of $X_H(x,v)$ in the coordinates $x^1,\ldots, x^n$ (and similarly for $X_V(x,v)$ below), this allows us to compute
 \begin{equation*}
\lim_{t\to 0^+} \frac{\partial \exp_x^{\varphi, i}}{\partial v^j}(tv) = \lim_{t\to 0^+} \frac{1}{t} \frac{\partial \gamma^i}{\partial v^j}(t,x,v) = \frac{\partial^2 \gamma^i}{\partial t \partial v^j}(0,x,v) = \frac{\partial^2 \gamma^i}{\partial v^j \partial t}(0,x,v) = \frac{\partial X_H^i}{\partial v^j}(x,v).
 \end{equation*}
The above limit is uniform in $v$ as $U_xM$ is compact and, together with \eqref{d0-psi-exp}, shows that $\exp_x^{\varphi}$ is of class $C^1$ at $0$.

Next, assume that $\varphi^t$ is an odd semi-spray flow, and that $\exp_x^{\varphi}$ is of class $C^2$ at $0$.
For any $v\in U_xM$, the curve $t\mapsto \exp_x^{\varphi}(tv)$ is a solution of \eqref{Dc'/dt=XV}, whose coordinate version reads
\begin{equation}\label{cov-coordinates}
 \frac{\partial^2\gamma^i}{\partial t^2}(t,x,v) + \Gamma_{jk}^i(\exp_x^{\varphi}(tv)) \frac{\partial\gamma^j}{\partial t}(t,x,v)\frac{\partial \gamma^k}{\partial t}(t,x,v) = X^i_V(\varphi^t(x,v)),
\end{equation}for all $i=1,\ldots, n$. Differentiating \eqref{gamma-i} with respect to $t$ twice, we obtain
\begin{equation}\label{d2gammai}
 \frac{\partial^2\gamma^i}{\partial t^2}(t,x,v) =
\frac{\partial^2 \exp_x^{\varphi,i}}{\partial v^j\partial v^k}(tv)v^jv^k.\end{equation}Using the $C^2$-regularity of $\exp_x^{\varphi}$ in order to take $t\to 0$ in \eqref{cov-coordinates}, together with \eqref{d2gammai} and the relations $(\partial \gamma^i/\partial t)(0,x,v) = v^i$ given by the semi-spray condition, it follows that
\begin{equation}\label{even-odd-again}
\left( \frac{\partial^2 \exp_x^{\varphi,i}}{\partial v^j\partial v^k}(0)+\Gamma^i_{jk}(x)\right)v^jv^k= X^i_V(x,v).
\end{equation}The left side of \eqref{even-odd-again} is manifestly even in $v$, while the right side is odd by assumption---this is only possible if both vanish individually. Therefore, $X_V(x,v)=0$ as claimed.
\end{proof}

\end{document}
